% Part: computability
% Chapter: recursive-functions
% Section: primitive-recursion

\documentclass[../../../include/open-logic-section]{subfiles}

\begin{document}

\olfileid[fr]{cmp}{rec}{pre}
\olsection{Récurrence primitive}

Tout entier naturel s'obtient à partir de~$0$ par un nombre fini
d'applications de l'opération successeur~$+1$ : c'est~$0$, ou le
successeur de~$0$, ou le successeur du successeur de~$0$, et ainsi
de suite. On peut tirer parti de cette propriété pour définir une
fonction $h\colon\Nat\to\Nat$ : (a)~on précise la valeur de $h$
en~$0$; (b)~on indique comment calculer $h(x+1)$ à partir de
$h(x)$. La première instruction donne directement $h(0)$ : la
fonction est donc définie en~$0$. La seconde, appliquée avec $x=0$,
permet de calculer $h(1)=h(0+1)$ à partir de~$h(0)$. Appliquée avec
$x=1$, elle donne $h(2)=h(1+1)$ à partir de~$h(1)$, et ainsi de
suite. Pour tout entier naturel~$x$, on parvient en un nombre fini
d'étapes à celle qui définit $h(x+1)$ à partir de~$h(x)$; la fonction
est donc définie sur tout~$\Nat$.\footnote{Correction éditoriale : la source inverse ici le sens du calcul en disant que $h(x)$ est défini à partir de $h(x+1)$. Les équations et la construction calculent au contraire $h(x+1)$ à partir de $h(x)$.}

Supposons par exemple que la fonction $h\colon\Nat\to\Nat$ soit
donnée par les deux équations suivantes :
\begin{align*}
h(0) & =  1\\
h(x+1) & =  2 \cdot h(x)
\end{align*}
Si nous savons déjà multiplier, ces équations fournissent les
informations requises en (a) et en~(b). En appliquant successivement
la seconde, nous obtenons :
\begin{align*}
  h(1) & = 2\cdot h(0) = 2,\\
  h(2) & = 2\cdot h(1) = 2\cdot 2,\\
  h(3) & = 2 \cdot h(2) = 2\cdot 2 \cdot 2,\\
  & \vdots
\end{align*}
La fonction ainsi définie est $h(x)=2^x$.

La propriété des entiers naturels rappelée plus haut garantit
l'unicité de la fonction~$h$ qui satisfait ces deux équations. Une
paire d'équations de ce type constitue une \emph{définition par
récurrence primitive} de~$h$. Elle est dite « récursive » parce que
la définition, et plus précisément sa seconde équation, fait
intervenir~$h$ elle-même dans le membre de droite. Elle est dite
« primitive » parce que, pour définir $h(x+1)$, elle n'utilise que
la valeur immédiatement précédente~$h(x)$. C'est le schéma le plus
simple de définition récursive d'une fonction sur~$\Nat$.

On peut aussi définir par récurrence primitive des fonctions plus
fondamentales, comme l'addition et la multiplication. Cette fois,
les fonctions ont deux arguments. On fixe l'un d'eux et l'on fait
porter la récurrence sur l'autre. Pour définir $\Add(x,y)$, par
exemple, on fixe~$x$ et l'on précise d'abord la valeur pour $y=0$,
puis celle pour $y+1$ à partir de celle pour~$y$. L'argument~$x$
étant fixé, il figure dans les deux membres des équations :
\begin{align*}
\Add(x,0) & =  x\\
\Add(x,y+1) & =  \Add(x,y)+1
\end{align*}
Ces équations déterminent $\Add(x,y)$ pour tous $x$ \emph{et}~$y$.
Pour calculer $\Add(2,3)$, par exemple, on les applique en fixant
$x=2$ : la première donne $\Add(2,0)=2$, puis la seconde donne
successivement $\Add(2,1)=2+1=3$, $\Add(2,2)=3+1=4$ et
$\Add(2,3)=4+1=5$.

Dans la seconde équation définissant $\Add$, le signe~$+$ du membre
de droite ne sert qu'à ajouter~$1$. Autrement dit, nous appliquons
la fonction successeur $\Succ(z)=z+1$ à la valeur précédente
$\Add(x,y)$ pour définir $\Add(x,y+1)$. On peut donc concevoir
cette définition récursive au moyen d'une seule fonction appliquée
à la valeur précédente. Il est toutefois sans inconvénient, et
parfois nécessaire, que cette fonction dépende aussi de $x$ et
de~$y$. Considérons :
\begin{align*}
  \Mult(x,0) & =  0 \\
  \Mult(x,y+1) & =  \Add(\Mult(x,y),x)
\end{align*}
Cette définition par récurrence primitive applique $\Add$ à la
valeur précédente $\Mult(x,y)$ et au premier argument~$x$. Elle
définit elle aussi $\Mult(x,y)$ pour tous $x$ et~$y$. On détermine
par exemple $\Mult(2,3)$ en calculant successivement $\Mult(2,0)$,
$\Mult(2,1)$, $\Mult(2,2)$ et~$\Mult(2,3)$ :
\begin{align*}
  \Mult(2,0) & = 0\\
  \Mult(2,1) & = \Mult(2,0+1) =
  \Add(\Mult(2,0), 2) = \Add(0, 2) = 2\\
  \Mult(2,2) & = \Mult(2,1+1) =
  \Add(\Mult(2,1), 2) = \Add(2, 2) = 4\\
  \Mult(2,3) & = \Mult(2,2+1) =
  \Add(\Mult(2,2), 2) = \Add(4, 2) = 6
\end{align*}

Voici le schéma général. Pour définir par récurrence primitive une
fonction $h(x_0,\dots,x_{k-1},y)$, on donne deux équations. La
première définit $h(x_0,\dots,x_{k-1},0)$ sans faire appel à~$h$.
La seconde définit $h(x_0,\dots,x_{k-1},y+1)$ à partir de la
valeur précédente $h(x_0,\dots,x_{k-1},y)$, des autres arguments
$x_0,\dots,x_{k-1}$ et de~$y$. Seule cette valeur immédiatement
précédente de~$h$ peut intervenir dans la seconde équation. Si les
opérations données par les membres de droite sont elles-mêmes des
fonctions $f$ et~$g$, le schéma prend la forme suivante :
\begin{align*}
  h(x_0, \dots, x_{k-1}, 0) & = f(x_0, \dots, x_{k-1})\\
  h(x_0, \dots, x_{k-1}, y+1) & =
  g(x_0, \dots, x_{k-1}, y, h(x_0, \dots, x_{k-1}, y))
\end{align*}
Pour $\Add$, on a $k=1$, $f(x_0)=x_0$ (la fonction identité) et
$g(x_0,y,z)=z+1$ (la fonction à trois arguments qui renvoie le
successeur du troisième) :
\begin{align*}
  \Add(x_0, 0) & = f(x_0) = x_0\\
  \Add(x_0, y+1) & = g(x_0, y, \Add(x_0, y)) =
  \Succ(\Add(x_0, y))
\end{align*}
Pour $\Mult$, on a $f(x_0)=0$ (la fonction constante de valeur~$0$)
et $g(x_0,y,z)=\Add(z,x_0)$ (la fonction à trois arguments qui
renvoie la somme du dernier et du premier) :
\begin{align*}
  \Mult(x_0,0) & =  f(x_0) = 0 \\
  \Mult(x_0,y+1) & =  g(x_0, y, \Mult(x_0,y)) =
  \Add(\Mult(x_0,y), x_0)
\end{align*}

\end{document}
